\documentclass[a4paper, 12pt]{article}

\usepackage[labelformat=simple]{subcaption}
\usepackage{graphicx}% LaTeX package to import graphics
\graphicspath{{images/}} % configuring the graphicx package



\usepackage{color}
\definecolor{darkblue}{rgb}{0,0,0.75}
\definecolor{darkgreen}{rgb}{0,0.75,0}
\usepackage[%pagebackref,
colorlinks, linkcolor=darkblue,
citecolor=darkgreen, urlcolor=blue, bookmarks=false, breaklinks]{hyperref}

\usepackage{cmap}

\usepackage[utf8]{inputenc}
\usepackage[T2A]{fontenc}
\usepackage[english,russian]{babel}

\usepackage{indentfirst}
\usepackage[left=2.5cm, right=1.5cm, top=2cm, bottom=2cm]{geometry}
\setlength\parindent{5ex}

\usepackage{fancyhdr}
\fancypagestyle{plain}{
	\fancyhf{}
	\chead{\thepage}
	\renewcommand{\headrulewidth}{0pt}}
\linespread{1.25}

\usepackage{tocloft}
\usepackage[ddmmyyyy]{datetime}

\usepackage{titlesec}
\titleformat{\section}[block]{\bfseries\large\centering}{\thesection}{1ex}{}
\titleformat{\subsection}[block]{\bfseries\normalsize\centering}{\thesubsection}{1ex}{}

\usepackage{amssymb, amsmath, amsthm, longtable, bbm}
\usepackage[matrix, arrow, curve]{xy}

\usepackage[geometry]{ifsym}

\renewcommand{\square}{\text{\SmallSquare}}

\binoppenalty=\maxdimen
\relpenalty=\maxdimen

\righthyphenmin=2
\sloppy

\newtheoremstyle{mytheoremstyle} % name
{\topsep}                    % Space above
{\topsep}                    % Space below
{\itshape}                   % Body font
{5ex}                           % Indent amount
{\bfseries}                   % Theorem head font
{.}                          % Punctuation after theorem head
{.5em}                       % Space after theorem head
{}  % Theorem head spec (can be left empty, meaning ‘normal’)

\newtheoremstyle{myremarkstyle} % name
{\topsep}                    % Space above
{\topsep}                    % Space below
{}                   % Body font
{5ex}                           % Indent amount
{\bfseries}                   % Theorem head font
{.}                          % Punctuation after theorem head
{.5em}                       % Space after theorem head
{}  % Theorem head spec (can be left empty, meaning ‘normal’)
\usepackage{titlesec}
\titleformat{\section}[block]{\normalfont\Large\bfseries}{\thesection}{1em}{}
\titlespacing*{\section}{0pt}{3ex}{2ex}

\theoremstyle{mytheoremstyle}
\newtheorem{theorem}{Теорема}[section]
\newtheorem{proposition}[theorem]{Предложение}
\newtheorem{corollary}[theorem]{Следствие}
\newtheorem{lemma}[theorem]{Лемма}
\newtheorem*{conjecture}{Гипотеза}

\theoremstyle{myremarkstyle}
\newtheorem{remark}[theorem]{Замечание}
\newtheorem{agreement}[theorem]{Соглашение}
\newtheorem{example}[theorem]{Пример}
\newtheorem*{examples}{Примеры}
\newtheorem*{questions}{Вопросы}
\newtheorem*{question}{Вопрос}
\newtheorem{definition}[theorem]{Определение}
\newtheorem{problem}[theorem]{Задача}

\numberwithin{equation}{section}

\makeatletter
\renewenvironment{proof}[1][\proofname]{\par\indent {\bfseries #1\@addpunct{.} }}{\qed}
\makeatother

\DeclareMathOperator{\Id}{Id}
\DeclareMathOperator{\Der}{Der}
\DeclareMathOperator{\id}{id}
\DeclareMathOperator{\chr}{char}
\DeclareMathOperator{\ad}{ad}
\DeclareMathOperator{\tr}{tr}
\DeclareMathOperator{\GL}{GL}
\DeclareMathOperator{\UT}{UT}
\DeclareMathOperator{\SL}{SL}
\DeclareMathOperator{\diff}{diff}
\DeclareMathOperator{\supp}{supp}
\DeclareMathOperator{\cosupp}{cosupp}
\DeclareMathOperator{\PGL}{PGL}
\DeclareMathOperator{\PSL}{PSL}
\DeclareMathOperator{\height}{ht}
\DeclareMathOperator{\End}{End}
\DeclareMathOperator{\Aut}{Aut}
\DeclareMathOperator{\diag}{diag}
\DeclareMathOperator{\rank}{rank}
\DeclareMathOperator{\length}{length}
\DeclareMathOperator{\sign}{sign}
\DeclareMathOperator{\Alt}{Alt}
\DeclareMathOperator{\Ann}{Ann}
\DeclareMathOperator{\Hom}{Hom}
\DeclareMathOperator{\op}{op}
\DeclareMathOperator{\PIexp}{PIexp}
\newcommand{\hatotimes}{\mathbin{\widehat{\otimes}}}
\DeclareMathOperator{\Ker}{Ker}
\DeclareMathOperator{\Coker}{Coker}
\renewcommand{\Im}{\mathop\mathrm{Im}}
\DeclareMathOperator{\Coim}{Coim}
\DeclareMathOperator{\Ext}{Ext}
\DeclareMathOperator{\Tor}{Tor}

%%%%%%%%%%%%%%% Fancy symbol \No %%%%%%%%%%%%%%%%%%%%%%%%%%%%%%%%%%%%%%%%%%%%%%

\DeclareRobustCommand{\No}{\ifmmode{\nfss@text{\textnumero}}\else\textnumero\fi} 


%%%%%%%%%%%%%%% Pullbacks van Tim van der Linden %%%%%%%%%%%%%%%%%%%%%%%%%%%%%

\newbox\skewpullbackbox
\setbox\skewpullbackbox=\hbox{\xy 0;<1mm,0mm>: \POS(4,0)\ar@{-} (-4,0) \ar@{-} (8,4)
	\endxy}
\newcommand{\skewpullback}{\copy\skewpullbackbox}

\newbox\skwepullbackbox
\setbox\skwepullbackbox=\hbox{\xy 0;<1mm,0mm>: \POS(16,0)\ar@{-} (10,0) \ar@{-} (12,4)
	\endxy}
\newcommand{\skwepullback}{\copy\skwepullbackbox}

\newbox\ksewpullbackbox
\setbox\ksewpullbackbox=\hbox{\xy 0;<1mm,0mm>: \POS(0,-8)\ar@{-} (0,-4) \ar@{-} (4,-4)
	\endxy}
\newcommand{\ksewpullback}{\copy\ksewpullbackbox}

\newbox\pullbackbox
\setbox\pullbackbox=\hbox{\xy 0;<1mm,0mm>: \POS(4,0)\ar@{-} (0,0) \ar@{-} (4,4)
	\endxy}
\newcommand{\pullback}{\copy\pullbackbox}

\newbox\pullbackabox
\setbox\pullbackabox=\hbox{\xy 0;<1mm,0mm>: \POS(-4,-6)\ar@{-} (-8,-6) \ar@{-} (-4,-2)
	\endxy}
\newcommand{\pullbacka}{\copy\pullbackabox}

\newbox\pullbackbbox
\setbox\pullbackbbox=\hbox{\xy 0;<1mm,0mm>: \POS(-4,-4)\ar@{-} (-8,-4) \ar@{-} (-4,0)
	\endxy}
\newcommand{\pullbackb}{\copy\pullbackbbox}

\newbox\pullbackcbox
\setbox\pullbackcbox=\hbox{\xy 0;<1mm,0mm>: \POS(4,-7)\ar@{-} (0,-7) \ar@{-} (4,-3)
	\endxy}
\newcommand{\pullbackc}{\copy\pullbackcbox}


\newbox\pullbackdbox
\setbox\pullbackdbox=\hbox{\xy 0;<1mm,0mm>:\POS(-10,-6)\ar@{-} (-14,-6) \ar@{-} (-10,-2)
	\endxy}
\newcommand{\pullbackd}{\copy\pullbackdbox}


\newbox\pushoutbox
\setbox\pushoutbox=\hbox{\xy 0;<1mm,0mm>: \POS(0,4)\ar@{-} (0,0) \ar@{-} (4,4)
	\endxy}
\newcommand{\pushout}{\copy\pushoutbox}

\newbox\pushoutabox
\setbox\pushoutabox=\hbox{\xy 0;<1mm,0mm>: \POS(2,8)\ar@{-} (2,4) \ar@{-} (8,8)
	\endxy}
\newcommand{\pushouta}{\copy\pushoutabox}

\addto\captionsrussian{% Replace "english" with the language you use
  \renewcommand{\contentsname}%
    {Содержание}%
}


\begin{document}

{\fontsize{14pt}{18}\selectfont
  
  \thispagestyle{empty}
  \begin{center}
		
		\vfill\vfill \ \\ {%\Large
			МОСКОВСКИЙ ГОСУДАРСТВЕННЫЙ УНИВЕРСИТЕТ \\
			имени М.В.~ЛОМОНОСОВА
			
			\medskip
			
			МЕХАНИКО-МАТЕМАТИЧЕСКИЙ ФАКУЛЬТЕТ


		}



{%\Large
	\vfill\vfill\vfill\vfill\vfill\vfill\vfill {%\Large
   
		Отчёт о задаче №4
		
	}
		
		
		\vfill{\Large
        
			\textbf{Удаление черных краев изображения} 
		}
		
			

			\vfill
			\begin{flushright}
				\begin{tabular}{r}
                \\
                \\
                \\
                \\
                \\
                \\
                \\
                \\
                
					Выполнила:\\
студентка 223 группы \\
Якшибаева Дарья Сергеевна
					
				\end{tabular}
			\end{flushright}
			
		}
		
		\vfill\vfill\vfill\vfill 
		
		Москва
		
		 2025
	\end{center}
}

    \newpage
    \tableofcontents
    \fontsize{14pt}{20pt}\selectfont
    \newpage
\pagestyle{plain}

\section{Описание задачи}

Задача состоит в написании программы на языке Python, которая обрезает черные края изображения. Ширина и высота изображения должны уменьшиться на размер обрезанных границ. В результате должно получиться изоображение без черных краев.

%\newpage
\section{Постановка задачи}

Входные данные:
\begin{itemize}
\setlength{\itemindent}{2cm}
Графический файл.\\
Допустимые форматы: JPEG, PNG, BMP
\end{itemize}

Выходные данные:
\begin{itemize}
\setlength{\itemindent}{1cm}
 Графический файл OutputFile (например, output.jpg) с обрезанными чёрными границами
\end{itemize}

Условия:
\begin{itemize}
\setlength{\itemindent}{2cm}
 Если изображение не содержит чёрных границ, размеры выходного файла совпадают с входным.\\
 Если все пиксели изображения чёрные, программа должна сохранить это же изображение без изменений.
\end{itemize}



\newpage
\section{Реализация задачи}

Под "чёрными краями" понимаются крайние строки/столбцы, полностью состоящие из пикселей с цветом (0, 0, 0-10).\\
Используется библиотека Pillow (PIL).Изображение конвертируется в RGB-формат.\\


Основные этапы
\begin{itemize}
\setlength{\itemindent}{2cm}
\item Загрузка входного изображения\\
 img = Image.open(sys.argv[1]) \\

\itemПреобразование в массив пикселей\\
pixels = np.array(img.convert('RGB'))\\
\itemОпределение границ обрезки\\
side = 0\\
    while side < x and np.all(pixels[side] <= 10):\\
    side += 1\\
  
    

\itemВыполнение обрезки\\
img.crop((left, top, right + 1, bottom + 1))\\
\itemСохранение результата\\
cropped.save(sys.argv[2])

\end{itemize}

\newpage
\section{Примеры выполнения}
\begin{figure}[h]
\centering
\begin{subfigure}{0.4\textwidth}
    \includegraphics[width=\linewidth]{input2.jpg}
    \caption{Исходное изображение}
\end{subfigure}
\hfill
\begin{subfigure}{0.4\textwidth}
    \includegraphics[width=\linewidth]{output2.jpg}
    \caption{Результат обработки}
\end{subfigure}
\caption{Пример работы программы}
\end{figure}

\begin{figure}[h]
\centering
\begin{subfigure}{0.4\textwidth}
    \includegraphics[width=\linewidth]{input1.jpg}
    \caption{Исходное изображение}
\end{subfigure}
\hfill
\begin{subfigure}{0.4\textwidth}
    \includegraphics[width=\linewidth]{output1.jpg}
    \caption{Результат обработки}
\end{subfigure}
\caption{Пример работы программы}
\end{figure}
\section{Вывод}
Программа успешно выполняет поставленную задачу по обработке изображений.





\end{document}